% ============================================================
% Tuyển tập bài tập Competitive Programming
% Dịch sang tiếng Việt — chuẩn HSG / ICPC
% ============================================================
\documentclass[12pt,a4paper]{article}

% ── Packages ─────────────────────────────────────────────────
\usepackage{tabularx}
\usepackage{booktabs}
\usepackage[utf8]{inputenc}
\usepackage[T5]{fontenc}
\usepackage[vietnamese]{babel}
\usepackage{amsmath,amssymb,amsthm}
\usepackage[margin=2.5cm]{geometry}
\usepackage{listings}
\usepackage{xcolor}
\usepackage{hyperref}
\usepackage{enumitem}
\usepackage{fancyhdr}
\usepackage{titlesec}
\usepackage{mdframed}
\usepackage{array}
\usepackage{parskip}

% ── Colors ───────────────────────────────────────────────────
\definecolor{sectioncolor}{RGB}{0,70,127}
\definecolor{codebg}{RGB}{245,245,245}
\definecolor{rulecolor}{RGB}{200,200,200}
\definecolor{framecolor}{RGB}{0,100,180}
\definecolor{inputbg}{RGB}{240,248,255}
\definecolor{constraintbg}{RGB}{255,250,235}

% ── Section styling ──────────────────────────────────────────
\titleformat{\section}
  {\normalfont\Large\bfseries\color{sectioncolor}}
  {\thesection}{1em}{}[{\color{sectioncolor}\titlerule[0.8pt]}]

\titleformat{\subsection}
  {\normalfont\large\bfseries\color{sectioncolor!80}}
  {\thesubsection}{1em}{}

% ── Header / Footer ──────────────────────────────────────────
\pagestyle{fancy}
\fancyhf{}
\fancyhead[L]{\textcolor{sectioncolor}{\small\textbf{Tuyển tập bài CP}}}
\fancyhead[R]{\textcolor{sectioncolor}{\small\thepage}}
\fancyfoot[C]{\textcolor{gray}{\small Dịch thuật: tiếng Việt — Chuẩn HSG/ICPC}}
\renewcommand{\headrulewidth}{0.4pt}

% ── Listings (code blocks) ────────────────────────────────────
\lstset{
  backgroundcolor=\color{codebg},
  basicstyle=\ttfamily\small,
  breaklines=true,
  frame=single,
  rulecolor=\color{rulecolor},
  framesep=4pt,
  xleftmargin=6pt,
  xrightmargin=6pt,
  aboveskip=6pt,
  belowskip=6pt,
}

% ── Macros ───────────────────────────────────────────────────
\newcommand{\problem}[2]{%
  \addcontentsline{toc}{section}{#1}%
  \begin{mdframed}[
    linecolor=framecolor,
    linewidth=1.5pt,
    backgroundcolor=white,
    topline=true, bottomline=true,
    leftline=true, rightline=true,
    innerleftmargin=10pt, innerrightmargin=10pt,
    innertopmargin=6pt, innerbottommargin=6pt,
  ]
  \textbf{\large #1} \hfill \textit{Nguồn: #2}
  \end{mdframed}
  \vspace{4pt}%
}

\newmdenv[
  linecolor=framecolor!60,
  linewidth=0.8pt,
  backgroundcolor=inputbg,
  innerleftmargin=8pt,
  innerrightmargin=8pt,
  innertopmargin=4pt,
  innerbottommargin=4pt,
  skipabove=4pt,
  skipbelow=4pt,
]{inputbox}

\newmdenv[
  linecolor=framecolor!60,
  linewidth=0.8pt,
  backgroundcolor=inputbg,
  innerleftmargin=8pt,
  innerrightmargin=8pt,
  innertopmargin=4pt,
  innerbottommargin=4pt,
  skipabove=4pt,
  skipbelow=4pt,
]{outputbox}

\newmdenv[
  linecolor=orange!70,
  linewidth=0.8pt,
  backgroundcolor=constraintbg,
  innerleftmargin=8pt,
  innerrightmargin=8pt,
  innertopmargin=4pt,
  innerbottommargin=4pt,
  skipabove=4pt,
  skipbelow=4pt,
]{constraintbox}

\newcommand{\inputformat}{%
  \vspace{6pt}\textbf{\textcolor{sectioncolor}{Dữ liệu vào}}\par\vspace{2pt}%
}
\newcommand{\outputformat}{%
  \vspace{6pt}\textbf{\textcolor{sectioncolor}{Dữ liệu ra}}\par\vspace{2pt}%
}
\newcommand{\constraints}{%
  \vspace{6pt}\textbf{\textcolor{sectioncolor}{Giới hạn}}\par\vspace{2pt}%
}
\newcommand{\example}{%
  \vspace{6pt}\textbf{\textcolor{sectioncolor}{Ví dụ}}\par\vspace{2pt}%
}
\newcommand{\explanation}{%
  \vspace{6pt}\textbf{\textcolor{sectioncolor}{Giải thích ví dụ}}\par\vspace{2pt}%
}

% ── Hyperref ─────────────────────────────────────────────────
\hypersetup{
  colorlinks=true,
  linkcolor=sectioncolor,
  urlcolor=sectioncolor!80,
  pdftitle={Tuyển tập bài CP},
  pdfauthor={Dịch thuật chuẩn HSG/ICPC},
}

% ── Miscellaneous ────────────────────────────────────────────
\setlength{\parindent}{0pt}
\setlength{\parskip}{6pt}

% ============================================================
\begin{document}
% ============================================================

% ── Title page ───────────────────────────────────────────────
\begin{titlepage}
  \centering
  \vspace*{3cm}
  {\Huge\bfseries\color{sectioncolor} Tuyển tập Bài tập\\[0.4em]
   Competitive Programming}\\[0.8em]
  {\large\color{gray} Cấu trúc Dữ liệu \& Giải thuật}\\[2em]
  \rule{\linewidth}{1pt}\\[1.5em]
  {\large Nguồn: Tập hợp từ nhiều nền tảng}\\[0.5em]
  {\large Ngôn ngữ: \textbf{Tiếng Việt} — Chuẩn HSG / ICPC}\\[2em]
  \rule{\linewidth}{1pt}\\[3em]
  {\normalsize\color{gray}
    Tài liệu dịch thuật phục vụ học tập và luyện tập lập trình thi đấu.\\
    Mọi nội dung được dịch sát nghĩa từ đề gốc, giữ nguyên cấu trúc logic.
  }
  \vfill
  {\small\color{gray} Ngày biên soạn: \today}
\end{titlepage}

% ── Table of contents ────────────────────────────────────────
\tableofcontents
\newpage

\problem{MEX Replacement on Tree}{Codeforces - 2219D}

\section*{Đề bài}
Cho một cây có gốc tại đỉnh $1$ gồm $n$ đỉnh và một hoán vị $p$ có độ dài $n$ gồm các số nguyên $0, 1, \ldots, n - 1$, trong đó $p_v$ là trọng số của đỉnh $v$.

Gọi $S_v$ là tập các trọng số của các đỉnh nằm trên đường đi từ $1$ đến $v$. Định nghĩa hàm $f$ là $f(v) = \text{MEX}(S_v)$, tức là $\text{MEX}$ của tập hợp gồm các trọng số của các đỉnh trên đường đi đơn từ gốc đến $v$ (bao gồm cả gốc và $v$). Hàm $\text{MEX}$ của một tập hợp các số nguyên $c_1, c_2, \ldots, c_k$ được định nghĩa là số nguyên không âm nhỏ nhất $x$ không xuất hiện trong tập $c$.

Hơn nữa, bạn có thể thực hiện thao tác sau:
\begin{itemize}
    \item Chọn một số nguyên $v$ ($1 \le v \le n$) và gán $f(v)$ cho $p_v$.
\end{itemize}

Mục tiêu của bạn là tìm giá trị lớn nhất của $\sum_{v = 1}^n f(v)$ sau khi thực hiện thao tác trên không quá một lần (có thể là không lần nào).

\begin{inputbox}
Mỗi test gồm nhiều test cases. Dòng đầu tiên chứa số lượng test cases $t$ ($1 \le t \le 2 \cdot 10^4$). Mô tả của các test cases như sau.

Dòng đầu tiên của mỗi test case chứa một số nguyên $n$ ($1 \le n \le 2 \cdot 10^5$) — số lượng đỉnh của cây.

Dòng thứ hai của mỗi test case chứa $n$ số nguyên $p_1, p_2, \ldots, p_n$ ($0 \le p_i \le n - 1$) — các phần tử của mảng $p$. Đảm bảo rằng mỗi số nguyên từ $0$ đến $n - 1$ xuất hiện chính xác một lần trong $p$.

$n - 1$ dòng tiếp theo, mỗi dòng chứa hai số nguyên $x_i$ và $y_i$ ($1 \le x_i, y_i \le n$), biểu diễn một cạnh giữa đỉnh $x_i$ và $y_i$. Đảm bảo rằng các cạnh này tạo thành một cây.

Đảm bảo rằng tổng của $n$ qua tất cả các test cases không vượt quá $2 \cdot 10^5$.
\end{inputbox}

\begin{outputbox}
Đối với mỗi test case, in ra một dòng chứa một số nguyên: giá trị lớn nhất của $\sum_{v = 1}^n f(v)$ sau khi thực hiện thao tác trên không quá một lần (có thể là không lần nào).
\end{outputbox}

\begin{constraintbox}
\begin{itemize}
    \item $1 \le t \le 2 \cdot 10^4$
    \item $1 \le n \le 2 \cdot 10^5$
    \item $0 \le p_i \le n - 1$
    \item $1 \le x_i, y_i \le n$
    \item Tổng của $n$ qua tất cả các test cases $\le 2 \cdot 10^5$
\end{itemize}
\end{constraintbox}

\section*{Ví dụ}

\textbf{Input}

\begin{lstlisting}
7
1
0
3
1 0 2
1 2
1 3
6
1 4 5 2 0 3
1 4
1 5
6 2
2 5
2 3
5
1 2 3 0 4
1 2
2 3
3 4
4 5
10
9 8 7 1 3 2 5 4 6 0
6 10
3 1
8 7
4 2
2 8
7 5
10 9
3 9
6 4
9
8 4 0 6 5 7 3 1 2
8 3
4 2
4 6
9 3
7 6
1 5
8 5
9 2
7
1 5 2 3 6 0 4
7 6
4 3
7 2
5 1
2 4
6 5
\end{lstlisting}


\textbf{Output}

\begin{lstlisting}
1
4
12
9
30
26
17
\end{lstlisting}


\explanation
Trong ví dụ thứ nhất, việc không thực hiện thao tác nào là tối ưu, do đó đáp án là $1$.

Trong ví dụ thứ hai, chúng ta có:
\begin{itemize}
    \item $f(1) = \text{MEX}(\{p_1\}) = \text{MEX}(\{1\}) = 0$
    \item $f(2) = \text{MEX}(\{p_1, p_2\}) = \text{MEX}(\{1, 0\}) = 2$
    \item $f(3) = \text{MEX}(\{p_1, p_3\}) = \text{MEX}(\{1, 2\}) = 0$
\end{itemize}

Và chúng ta có các lựa chọn sau:
\begin{enumerate}
    \item Không thực hiện thao tác nào, tổng $\sum_{v = 1}^n f(v) = 2$.
    \item Thực hiện thao tác trên đỉnh $1$ sẽ gán $f(1) = 0$ cho $p_1$, làm thay đổi $f(1), f(2), f(3)$ thành $1$, dẫn đến $\sum_{v = 1}^n f(v) = 3$.
    \item Thực hiện thao tác trên đỉnh $2$ sẽ gán $f(2) = 2$ cho $p_2$, làm thay đổi $f(2)$ thành $0$, dẫn đến $\sum_{v = 1}^n f(v) = 0$.
    \item Thực hiện thao tác trên đỉnh $3$ sẽ gán $f(3) = 0$ cho $p_3$, làm thay đổi $f(3)$ thành $2$, dẫn đến $\sum_{v = 1}^n f(v) = 4$.
\end{enumerate}

Vậy đáp án lớn nhất là $4$, thu được bằng cách thực hiện thao tác trên đỉnh $3$.

\newpage
\problem{C. Square Subsets}{Codeforces - 895C}

\section*{Đề bài}
Petya lại đi học muộn. Thầy giáo đã giao cho cậu ấy một bài tập thêm. Cho một mảng $a$, Petya cần tìm số cách khác nhau để chọn ra một tập con khác rỗng gồm các phần tử từ mảng sao cho tích của chúng bằng với bình phương của một số nguyên nào đó.

Hai cách chọn được coi là khác nhau nếu tập hợp các chỉ số của các phần tử được chọn ở những cách này là khác nhau.

Vì đáp số có thể rất lớn, bạn nên tìm đáp án theo modulo $10^9+7$.

\section*{Dữ liệu vào}
- Dòng đầu tiên chứa một số nguyên $n$ ($1 \le n \le 10^5$) — số lượng phần tử của mảng.
- Dòng thứ hai chứa $n$ số nguyên $a_i$ ($1 \le a_i \le 70$) — các phần tử của mảng.

\section*{Kết quả}
- In ra một số nguyên duy nhất — số cách khác nhau để chọn một số phần tử sao cho tích của chúng là một số chính phương, lấy modulo $10^9+7$.

\section*{Ví dụ}


\begin{inputbox}4
1 1 1 1
\end{inputbox}
\begin{outputbox}15
\end{outputbox}


\begin{inputbox}4
2 2 2 2
\end{inputbox}
\begin{outputbox}7
\end{outputbox}


\begin{inputbox}5
1 2 4 5 8
\end{inputbox}
\begin{outputbox}7
\end{outputbox}



\section*{Giải thích}
\begin{itemize}
    \item \textbf{Ví dụ 1:} Mảng $a = [1, 1, 1, 1]$. Tích của bất kỳ tập con khác rỗng nào cũng đều bằng $1 = 1^2$.
Số lượng tập con khác rỗng có thể chọn là $2^4 - 1 = 16 - 1 = 15$.
Do đó, có $15$ cách chọn thoả mãn.

    \item \textbf{Ví dụ 2:} Mảng $a = [2, 2, 2, 2]$. Để tích là một số chính phương, ta cần chọn một số lượng chẵn các số $2$.
    \begin{itemize}
        \item Số cách chọn 2 phần tử $2$: $\binom{4}{2} = 6$. Tích là $4 = 2^2$.
        \item Số cách chọn 4 phần tử $2$: $\binom{4}{4} = 1$. Tích là $16 = 4^2$.
    \end{itemize}
    Tổng cộng có $6 + 1 = 7$ cách.

    \item \textbf{Ví dụ 3:} Mảng $a = [1, 2, 4, 5, 8]$. Ta có thể phân tích các phần tử ra thừa số nguyên tố:
    \begin{itemize}
        \item $1 = 1$
        \item $2 = 2^1$
        \item $4 = 2^2$
        \item $5 = 5^1$
        \item $8 = 2^3$
    \end{itemize}
    Để tích là một số chính phương, số mũ của các thừa số nguyên tố trong tích phải là số chẵn.
    Các tập con có tích là số chính phương bao gồm:
    \begin{itemize}
        \item $\{1\}$ (tích $1 = 1^2$)
        \item $\{4\}$ (tích $4 = 2^2$)
        \item $\{1, 4\}$ (tích $4 = 2^2$)
        \item $\{2, 8\}$ (tích $16 = 4^2$)
        \item $\{1, 2, 8\}$ (tích $16 = 4^2$)
        \item $\{4, 2, 8\}$ (tích $64 = 8^2$)
        \item $\{1, 4, 2, 8\}$ (tích $64 = 8^2$)
    \end{itemize}
    Tổng cộng có $7$ tập con thỏa mãn. (Phần tử $5$ có bậc lẻ nên bất kỳ tập con nào chứa $5$ sẽ không thể có tích là số chính phương).
\end{itemize}

\newpage
\problem{F. The Cake Is a Lie}{Codeforces - 2232F}

\section*{Đề bài}
Để kết thúc bữa tiệc, Alice quyết định nấu một số chiếc bánh kếp cho bạn bè của mình.

Alice có $n$ chiếc bánh kếp chưa nấu chín nhưng chỉ có $2$ cái chảo. Ban đầu, $\min(n, 2)$ chiếc bánh kếp đầu tiên được đặt trên chảo và tất cả các chiếc bánh kếp đều có độ chín bằng $0$.

Để nấu bánh kếp, Alice có thể thực hiện $2$ thao tác sau với số lần tùy ý:
\begin{itemize}
    \item Nấu hai chiếc bánh trên chảo trong $1$ phút. Độ chín của chiếc bánh trên chảo thứ nhất tăng thêm $a$, và độ chín của chiếc bánh trên chảo thứ hai tăng thêm $b$.
    \item Dọn (lấy ra) chiếc bánh thứ nhất, di chuyển chiếc bánh thứ hai sang chảo thứ nhất, và thêm một chiếc bánh kếp mới chưa nấu vào chảo thứ hai (nếu còn chiếc bánh nào hoàn toàn chưa nấu). Lưu ý rằng thao tác này không tốn thời gian, và có thể được thực hiện nhiều lần liên tiếp.
\end{itemize}

Lưu ý rằng nếu chỉ còn lại một chiếc bánh, nó sẽ nằm ở chảo thứ nhất, nhận thêm $a$ độ chín mỗi phút cho đến khi Alice quyết định dọn nó.

Một chiếc bánh kếp được nấu chín hoàn hảo khi và chỉ khi độ chín của nó bằng chính xác $k$.

Giúp Alice tìm số lượng bánh kếp chín hoàn hảo tối đa mà cô ấy có thể dọn.

\section*{Dữ liệu vào}
Mỗi test bao gồm nhiều test case. Dòng đầu tiên chứa số lượng test case $t$ ($1 \le t \le 1000$). Phần mô tả các test case theo sau.

Dòng duy nhất trong mỗi test case chứa 4 số nguyên $n$, $a$, $b$, và $k$ ($1 \le n, a, b, k \le 10^9$) - số lượng bánh kếp Alice có, lượng độ chín chiếc bánh ở chảo thứ nhất tăng thêm sau mỗi phút, lượng độ chín chiếc bánh ở chảo thứ hai tăng thêm sau mỗi phút, và lượng độ chín cần thiết để một chiếc bánh được nấu chín hoàn hảo.

\section*{Dữ liệu ra}
Với mỗi test case, in ra số lượng bánh kếp chín hoàn hảo tối đa mà Alice có thể làm.

\section*{Ví dụ}

\begin{inputbox}
7
17 1 1 1
123456789 987 654 321
3 11 37 111111
987654321 1 2 123456789
100 1 2 1
100 1 2 2
1 1 1 1
\end{inputbox}
\begin{outputbox}
17
0
2
987654320
50
67
1
\end{outputbox}


\explanation
Gọi $x_i$ là số phút chiếc bánh kếp thứ $i$ được nấu trên chảo thứ nhất.
Khi đó, độ chín của chiếc bánh thứ $i$ (với $i \ge 2$) sẽ là:
\[ C_i = b \cdot x_{i-1} + a \cdot x_i \]
Riêng chiếc bánh đầu tiên: $C_1 = a \cdot x_1$.

Chúng ta muốn tối đa hóa số lượng bánh kếp có $C_i = k$.
Mỗi chiếc bánh đạt độ chín hoàn hảo tương đương với một phương trình:
$b \cdot x_{i-1} + a \cdot x_i = k$ (với điều kiện $x_{i-1}, x_i \ge 0$).

Nếu xem mỗi giá trị $x \ge 0$ là một đỉnh của đồ thị, và có cung có hướng $u \to v$ nếu $b \cdot u + a \cdot v = k$, đồ thị này sẽ bao gồm các đường đi rời rạc hoặc các chu trình, vì mỗi đỉnh có bậc vào và bậc ra tối đa là 1. Một chuỗi các chiếc bánh liên tiếp đạt độ chín hoàn hảo sẽ tương ứng với một đường đi trên đồ thị này. 

Giả sử chúng ta tìm được đường đi dài nhất trên đồ thị có số đỉnh là $L_{max}$. Mỗi khi sử dụng đường đi này, chúng ta tiêu tốn $L_{max}$ chiếc bánh kếp và thu được $L_{max} - 1$ chiếc bánh hoàn hảo. Riêng chiếc bánh đầu tiên có thể cho thêm 1 chiếc bánh hoàn hảo nếu chúng ta bắt đầu tại $x_1 = k/a$ (khi $k$ chia hết cho $a$). Gọi $L_0$ là độ dài của đường đi bắt đầu từ đỉnh $k/a$ ($L_0 = 0$ nếu $k$ không chia hết cho $a$).

Công thức tính số lượng bánh hoàn hảo tối đa là:
\[ \text{Ans} = n - \left\lceil \frac{\max(0, n - L_0)}{L_{max}} \right\rceil \]
(Lưu ý: Nếu $L_{max} = \infty$, nghĩa là có chu trình, $\text{Ans} = n$ nếu $L_0 = \infty$, ngược lại $\text{Ans} = n - 1$).

Đặc điểm của đồ thị trong bài:
\begin{itemize}
    \item Nếu $a = b$: Nếu $k$ chia hết cho $a$, ta có chu trình $0 \leftrightarrow k/a$, do đó $L_{max} = \infty$ và $L_0 = \infty$. Ngược lại, đồ thị không có cung nào, $L_{max} = 0, L_0 = 0$.
    \item Nếu $a \neq b$: Đồ thị không có chu trình. Độ dài đường đi lớn nhất $L_{max}$ rất nhỏ ($L_{max} \le \log k \le 40$). Ta có thể tính $L_0$ bằng cách mô phỏng trực tiếp từ đỉnh $k/a$. Để tìm $L_{max}$, ta có thể thử duyệt qua mọi độ dài $L \ge 1$ để kiểm tra sự tồn tại đường đi hợp lệ.
\end{itemize}

\textbf{Phân tích từng bước cho các Test Cases:}
\begin{itemize}
    \item \textbf{Test case 1:} $n=17, a=1, b=1, k=1$. Có $a=b=1$ và $k$ chia hết cho $a$. Chu trình tạo ra là $x_i = 1, 0, 1, 0 \dots$ lặp lại vô hạn $\implies L_0 = \infty, L_{max} = \infty$. Kết quả: $17 - \lceil \max(0, 17-\infty)/\infty \rceil = 17 - 0 = 17$.
    \item \textbf{Test case 2:} $n=123456789, a=987, b=654, k=321$. Phương trình $654 x_{i-1} + 987 x_i = 321$ vô nghiệm với số tự nhiên $\implies L_0 = 0, L_{max} = 0$. Kết quả: 0.
    \item \textbf{Test case 3:} $n=3, a=11, b=37, k=111111$. Tính $L_0$: bắt đầu tại $k/a = 10101$, đỉnh tiếp theo là $(111111 - 37 \times 10101) / 11 < 0 \implies L_0 = 1$. Đường đi dài nhất trên đồ thị là $3003 \to 0 \to 10101 \implies L_{max} = 3$. Kết quả: $3 - \lceil (3 - 1)/3 \rceil = 2$.
    \item \textbf{Test case 4:} $n=987654321, a=1, b=2, k=123456789$. $k$ chia hết cho $a+b = 3$, điểm bất động $x^* = 123456789 / 3 = 41152263$. Đồ thị có một tự vòng (chu trình) tại đỉnh $x^* \implies L_{max} = \infty$. Bắt đầu từ $k/a$ có $L_0 = 1$. Kết quả: $n - \lceil (n-1)/\infty \rceil = n - 1 = 987654320$.
    \item \textbf{Test case 5:} $n=100, a=1, b=2, k=1$. Bắt đầu từ $k/a = 1 \implies L_0 = 1$. Cạnh hợp lệ duy nhất là $0 \to 1 \implies L_{max} = 2$. Kết quả: $100 - \lceil (100 - 1)/2 \rceil = 50$.
    \item \textbf{Test case 6:} $n=100, a=1, b=2, k=2$. Bắt đầu từ $k/a = 2 \implies L_0 = 1$. Đường đi dài nhất là $1 \to 0 \to 2 \implies L_{max} = 3$. Kết quả: $100 - \lceil 99/3 \rceil = 67$.
\end{itemize}

\newpage
\problem{E. Snaking Arrangement}{Codeforces - 2232E}

\section*{Đề bài}
Đã đến lúc cho ngôi sao của buổi tiệc: một chiếc bánh kem khổng lồ có kích thước $n \times n$.

Alice và các bạn muốn trang trí chiếc bánh bằng kem tươi có hình dạng những con rắn. Một con rắn có kích thước $k$ được định nghĩa là một đường đi có độ dài $k$ bắt đầu từ một ô của chiếc bánh và chỉ đi xuống dưới hoặc sang phải.

Nhờ việc nghiên cứu kỹ lưỡng về bánh kem, Alice biết rằng cách tốt nhất để trang trí chiếc bánh này là đặt $n$ con rắn sao cho con rắn thứ $i$ có kích thước là $2 \cdot i - 1$. Cô ấy đã chuẩn bị sẵn tất cả kem tươi; tuy nhiên, một vài người bạn của cô quá hào hứng và đã bắt đầu đặt vài con rắn lên chiếc bánh trước khi Alice kịp lên kế hoạch trang trí cụ thể.

May mắn thay, Alice nhận ra rằng việc trang trí chiếc bánh vẫn có thể thực hiện được. Hãy giúp cô ấy tìm số cách trang trí chiếc bánh mà không cần di chuyển bất kỳ con rắn nào các bạn của cô đã đặt. Vì số cách có thể rất lớn, hãy in ra đáp án theo modulo $10^9 + 7$.

Hai cấu hình được coi là khác nhau nếu có ít nhất một ô vuông bị chiếm bởi một con rắn khác.

\section*{Dữ liệu vào}
Mỗi test bao gồm nhiều test cases. Dòng đầu tiên chứa số lượng test cases $t$ ($1 \le t \le 1000$). Mô tả của các test cases như sau.

Dòng đầu tiên của mỗi test case chứa hai số nguyên $n$ và $k$ ($1 \le n \le 5000$, $0 \le k \le n$) — kích thước của chiếc bánh và số lượng con rắn mà các bạn của cô đã đặt.

$2 \cdot k$ dòng tiếp theo chứa thông tin về $k$ con rắn đã được đặt. Mỗi cặp dòng liên tiếp chứa thông tin về một con rắn được đặt theo định dạng sau:
\begin{itemize}
    \item Dòng thứ nhất chứa số nguyên $s$ ($1 \le s \le 2n-1$, $s$ là số lẻ) — kích thước của con rắn.
    \item Dòng thứ hai chứa $r$, $c$ ($1 \le r, c \le n$) — hàng và cột bắt đầu của con rắn. Tiếp theo là một xâu độ dài $s-1$ chỉ chứa các chữ cái \texttt{R} và \texttt{D}, trong đó \texttt{R} có nghĩa là ô tiếp theo nằm bên phải ô hiện tại, và \texttt{D} có nghĩa là ô tiếp theo nằm bên dưới ô hiện tại. Lưu ý rằng nếu $s = 1$, xâu này sẽ rỗng.
\end{itemize}

Đảm bảo rằng không có hai con rắn nào đè lên nhau, không có hai con rắn nào có cùng kích thước và mọi con rắn đều nằm gọn trong phạm vi của chiếc bánh.

Đảm bảo rằng đáp án cho mỗi test case luôn lớn hơn không.

Đảm bảo rằng tổng của $n$ qua tất cả các test cases không vượt quá $5000$, và tổng của $s$ qua tất cả các test cases không vượt quá $4 \cdot 10^5$.

\section*{Kết quả}
Đối với mỗi test case, in ra một số nguyên: số cách mà Alice có thể trang trí chiếc bánh theo modulo $10^9 + 7$.

\section*{Ví dụ}

\begin{inputbox}
5
3 1
5
1 1 RDRD
4 1
1
3 2
3 1
5
1 1 RRDD
3 1
3
2 1 DR
4567 0
\end{inputbox}
\begin{outputbox}
1
6
2
2
833729690
\end{outputbox}


\section*{Giải thích}
\begin{itemize}
    \item \textbf{Test case 1:} $n = 3, k = 1$. Có $n=3$ con rắn cần đặt với kích thước lần lượt là $1, 3, 5$. Con rắn đã được đặt có kích thước $s=5$, xuất phát tại $(1, 1)$ và di chuyển theo đường \texttt{RDRD} qua các ô $(1, 1) \to (1, 2) \to (2, 2) \to (2, 3) \to (3, 3)$. Các ô còn trống là $(1, 3), (2, 1), (3, 1), (3, 2)$. Con rắn kích thước 3 bắt buộc phải điền vào $(2, 1) \to (3, 1) \to (3, 2)$. Con rắn kích thước 1 điền vào ô $(1, 3)$. Chỉ có duy nhất $1$ cách sắp xếp.
    \item \textbf{Test case 2:} $n = 4, k = 1$. Con rắn đã đặt có kích thước $s=1$ nằm ở vị trí $(3, 2)$. Các con rắn cần đặt có kích thước $3, 5, 7$. Có tổng cộng $6$ cách hợp lệ để xếp các con rắn này vào phần còn lại của chiếc bánh.
    \item \textbf{Test case 3:} $n = 3, k = 1$. Con rắn kích thước $5$ bắt đầu tại $(1, 1)$ và đi theo đường \texttt{RRDD} qua $(1, 1) \to (1, 2) \to (1, 3) \to (2, 3) \to (3, 3)$. Các ô trống còn lại tạo thành hình vuông $2 \times 2$ ở góc dưới bên trái gồm các ô $(2, 1), (2, 2), (3, 1), (3, 2)$. Con rắn kích thước $3$ có thể đi $(2, 1) \to (3, 1) \to (3, 2)$ và con rắn kích thước $1$ nằm ở $(2, 2)$, hoặc con rắn kích thước $3$ đi $(2, 1) \to (2, 2) \to (3, 2)$ và kích thước $1$ nằm ở $(3, 1)$. Có $2$ cách sắp xếp.
    \item \textbf{Test case 4:} $n = 3, k = 1$. Con rắn kích thước $3$ xuất phát tại $(2, 1)$, đi theo đường \texttt{DR} qua $(2, 1) \to (3, 1) \to (3, 2)$. Các con rắn cần đặt có kích thước $1$ và $5$. Con rắn kích thước $5$ có độ dài lớn nhất nên bắt buộc phải nằm trên đường đi từ $(1, 1)$ đến $(3, 3)$. Nó có thể đi theo đường $(1, 1) \to (1, 2) \to (1, 3) \to (2, 3) \to (3, 3)$ (để lại $(2, 2)$ cho rắn kích thước $1$), hoặc $(1, 1) \to (1, 2) \to (2, 2) \to (2, 3) \to (3, 3)$ (để lại $(1, 3)$ cho rắn kích thước $1$). Có $2$ cách sắp xếp.
    \item \textbf{Test case 5:} Chiếc bánh có kích thước lớn ($n = 4567$) và không có con rắn nào được đặt trước. Số lượng cách trang trí chiếc bánh là rất lớn, và kết quả được in ra là phần dư sau khi chia cho $10^9 + 7$.
\end{itemize}

\newpage
\problem{D. Chiếc Bánh Tầng Ma Thuật}{Codeforces - 2232D}

\section*{Đề bài}
Alice đã hoàn thành một chiếc bánh tầng ma thuật với $n$ tầng ma thuật, trong đó mỗi tầng nhỏ hơn tất cả các tầng bên dưới nó (tức là tầng thứ nhất là nhỏ nhất, tầng thứ $n$ là lớn nhất). Bây giờ, cô ấy cần bạn giúp vận chuyển nó từ nhà bếp của cô ấy đến địa điểm tổ chức tiệc. Vì không thể di chuyển toàn bộ chiếc bánh cùng một lúc, cô ấy cũng chuẩn bị một nhà kho cho bạn để tạo điều kiện thuận lợi cho quá trình vận chuyển.

Vì chiếc bánh là ma thuật, tầng bánh thứ $i$ có thể di chuyển được khi và chỉ khi có chính xác $a_i$ tầng bánh ở phía trên nó.

Trong mỗi lần di chuyển, bạn có thể chọn \textbf{chính xác một} tầng bánh có thể di chuyển được từ bất kỳ vị trí nào và di chuyển nó lên trên chiếc bánh tầng ở bất kỳ vị trí nào khác. Tuy nhiên, để bảo tồn cấu trúc của chiếc bánh, tầng được di chuyển phải nằm trên tầng lớn hơn nó một cách nghiêm ngặt nếu có một chiếc bánh tầng ở điểm đến. Ví dụ, bạn không thể di chuyển một tầng có kích thước $4$ lên trên một vị trí đã có một tầng có kích thước $3$.

Bữa tiệc sắp bắt đầu, vì vậy chúng ta cần di chuyển nhanh chóng. Hãy giúp Alice vận chuyển chiếc bánh đến bữa tiệc trong vòng $2^n$ lần di chuyển hoặc thông báo rằng điều đó là không thể.

\section*{Dữ liệu vào}
Mỗi bài kiểm tra chứa nhiều test case. Dòng đầu tiên chứa số lượng test case $t$ ($1 \le t \le 10000$). Mô tả của các test case theo sau.

Dòng đầu tiên của mỗi test case chứa một số nguyên $n$ ($1 \le n \le 20$) — số tầng trong chiếc bánh tầng ma thuật.

Dòng thứ hai của mỗi test case chứa $n$ số nguyên $a_1, a_2, \ldots, a_n$ ($0 \le a_i \le n$), biểu thị có bao nhiêu tầng bánh cần phải ở trên tầng thứ $i$ để nó có thể di chuyển được.

Đảm bảo rằng tổng của $2^n$ trên tất cả các test case không vượt quá $2^{20}$.

\section*{Dữ liệu ra}
Đối với mỗi test case, nếu không thể di chuyển chiếc bánh tầng ma thuật đến địa điểm tổ chức tiệc, in ra \texttt{NO}.

Nếu không, trên dòng đầu tiên, in ra \texttt{YES}.

Sau đó, trên dòng tiếp theo, in ra một số nguyên $k$ ($0 \le k \le 2^n$) — số lần di chuyển bạn thực hiện để vận chuyển chiếc bánh.

Trong $k$ dòng tiếp theo, in ra các số nguyên $\mathrm{id}$, $\mathrm{from}$, $\mathrm{to}$ ($1 \le \mathrm{id} \le n$, $1 \le \mathrm{from}, \mathrm{to} \le 3$): tầng bánh bạn muốn di chuyển, vị trí hiện tại của nó, và điểm đến của nó, trong đó $1$ đại diện cho nhà bếp của Alice, $2$ đại diện cho nhà kho của Alice, và $3$ đại diện cho địa điểm tổ chức tiệc.

\section*{Ví dụ}

\begin{inputbox}
3
3
0 0 0
3
0 1 2
3
2 2 2
\end{inputbox}\begin{outputbox}
YES
7
1 1 3
2 1 2
1 3 2
3 1 3
1 2 1
2 2 3
1 1 3
YES
3
3 1 3
2 1 3
1 1 3
NO
\end{outputbox}%


\section*{Giải thích}
\begin{itemize}
    \item \textbf{Test case 1:} $n = 3$, $a = [0, 0, 0]$. Mỗi tầng yêu cầu $0$ tầng ở trên mới di chuyển được (nghĩa là nó phải nằm trên cùng). Đây chính là bài toán Tháp Hà Nội cổ điển. Các thao tác di chuyển lần lượt như sau (tốn $2^3 - 1 = 7$ bước):
    \begin{enumerate}
        \item Chuyển tầng 1 từ nhà bếp (1) sang địa điểm tiệc (3).
        \item Chuyển tầng 2 từ nhà bếp (1) sang nhà kho (2).
        \item Chuyển tầng 1 từ địa điểm tiệc (3) sang nhà kho (2).
        \item Chuyển tầng 3 từ nhà bếp (1) sang địa điểm tiệc (3).
        \item Chuyển tầng 1 từ nhà kho (2) sang nhà bếp (1).
        \item Chuyển tầng 2 từ nhà kho (2) sang địa điểm tiệc (3).
        \item Chuyển tầng 1 từ nhà bếp (1) sang địa điểm tiệc (3).
    \end{enumerate}

    \item \textbf{Test case 2:} $n = 3$, $a = [0, 1, 2]$. Ban đầu ở vị trí 1, chiếc bánh có tầng 1 trên cùng, tầng 2 ở giữa, tầng 3 dưới cùng.
    \begin{itemize}
        \item Tầng 3 có chính xác 2 tầng ở trên nó (tầng 1 và 2), thoả mãn $a_3=2$, nên tầng 3 có thể di chuyển ngay. Chuyển tầng 3 từ nhà bếp (1) sang địa điểm tiệc (3). Kích thước 3 là lớn nhất, đặt tại vị trí trống là hợp lệ.
        \item Ở vị trí 1 còn lại tầng 1 trên tầng 2. Tầng 2 có chính xác 1 tầng ở trên nó, thoả mãn $a_2=1$. Chuyển tầng 2 từ nhà bếp (1) sang địa điểm tiệc (3). Đặt tầng 2 lên tầng 3 là hợp lệ.
        \item Ở vị trí 1 còn lại tầng 1. Tầng 1 có 0 tầng ở trên, thoả mãn $a_1=0$. Chuyển tầng 1 từ nhà bếp (1) sang địa điểm tiệc (3). Đặt tầng 1 lên tầng 2 là hợp lệ.
    \end{itemize}
    Toàn bộ quá trình tốn 3 bước di chuyển.

    \item \textbf{Test case 3:} $n = 3$, $a = [2, 2, 2]$. Ở trạng thái ban đầu, chỉ có tầng 3 là có đủ 2 tầng ở trên và có thể di chuyển. Nếu ta di chuyển tầng 3, ở vị trí 1 sẽ chỉ còn 2 tầng (tầng 1 và 2), cả hai đều không thể thỏa mãn điều kiện có 2 tầng ở trên nó. Do đó, không thể vận chuyển toàn bộ chiếc bánh, và kết quả là \texttt{NO}.
\end{itemize}

\newpage
\problem{C2. Sắp Xếp Chỗ Ngồi (Bản Khó)}{Codeforces - 2232C2}

\section*{Đề bài}
Đây là phiên bản Khó của bài toán. Sự khác biệt giữa các phiên bản là ở phiên bản này, giới hạn của $n$, $x$, $s$, $t$ lớn hơn. Bạn chỉ có thể hack nếu bạn đã giải được tất cả các phiên bản của bài toán này.

Bạn bè của Alice đã đến bữa tiệc, và bây giờ họ đang xếp hàng để vào bữa tiệc.

Có $x$ cái bàn tại bữa tiệc, mỗi bàn có $s$ chỗ ngồi. Mỗi chỗ ngồi chỉ có thể chứa một người.

Mỗi người bạn có một trong ba tính cách sau:
\begin{itemize}
    \item Người hướng nội (\texttt{I}) phải ngồi ở một bàn trống.
    \item Người hướng ngoại (\texttt{E}) phải ngồi ở một bàn không trống.
    \item Người hướng trung (\texttt{A}) có thể ngồi ở bất kỳ bàn nào.
\end{itemize}

Ban đầu, mọi chỗ ngồi đều trống. Tuy nhiên, do Alice mải ăn bánh, bạn bè của cô ấy đã tự động xếp thành một hàng và Alice không thể thay đổi thứ tự của họ. Đối với mỗi người trong hàng, Alice phải xếp họ vào một bàn hoặc đuổi họ ra khỏi bữa tiệc. Mỗi người được xếp chỗ trước khi người tiếp theo được phân bàn.

Muốn có thật nhiều niềm vui tại bữa tiệc, Alice cần xếp chỗ cho càng nhiều người càng tốt. Hãy giúp cô ấy tìm số lượng bạn bè tối đa mà cô ấy có thể xếp chỗ tại bữa tiệc.

Lưu ý rằng một khi một người bạn đã được xếp chỗ, họ không được phép di chuyển ngay cả khi họ không còn ngồi theo đúng tính cách của họ nữa.

\section*{Dữ liệu vào}
Mỗi test bao gồm nhiều test cases. Dòng đầu tiên chứa số lượng test cases $t$ ($1 \le t \le 10^4$). Theo sau là mô tả của các test cases.

Dòng đầu tiên của mỗi test case chứa ba số nguyên $n$, $x$, và $s$ ($1 \le n, x, s \le 2\cdot 10^5$) — lần lượt là số lượng bạn bè của Alice, số lượng bàn, và số chỗ ngồi mỗi bàn. 

Dòng thứ hai chứa một chuỗi $u$ độ dài $n$ chỉ bao gồm các chữ cái \texttt{A}, \texttt{E}, và \texttt{I}, lần lượt đại diện cho người hướng trung, người hướng ngoại và người hướng nội.

Dữ liệu đảm bảo tổng của $n$ trên tất cả các test cases không vượt quá $2 \cdot 10^5$.

\section*{Dữ liệu ra}
Đối với mỗi test case, in ra một số nguyên: số người tối đa được xếp chỗ.

\section*{Ví dụ}


\begin{inputbox}6
5 2 2
EIAIE
20 5 5
AEIEEEEIEAAEIEEEEIEA
8 2 4
AAAAAIEE
8 4 2
AIEAEAAI
8 3 3
AIEAEAAI
4 2 2
IAEE
\end{inputbox}
\begin{outputbox}4
20
7
7
7
4
\end{outputbox}



\section*{Giải thích}
Trong test case đầu tiên, có $2$ bàn, mỗi bàn có $2$ chỗ ngồi. Dưới đây là một trong những cách để đạt được số người ngồi tối đa:
\begin{itemize}
    \item Người thứ nhất là người hướng ngoại. Vì tất cả các bàn đều trống, họ đành phải rời bữa tiệc.
    \item Người thứ hai là người hướng nội. Alice có thể xếp họ vào bàn thứ nhất, hiện đang trống.
    \item Người thứ ba là người hướng trung. Alice có thể xếp họ vào bàn thứ nhất.
    \item Người thứ tư là người hướng nội. Alice có thể xếp họ vào bàn thứ hai, hiện đang trống.
    \item Người thứ năm là người hướng ngoại. Alice có thể xếp họ vào bàn thứ hai, hiện không trống.
\end{itemize}
Như vậy, có $4$ người được ngồi tại bữa tiệc. Con số này là tối đa vì bữa tiệc chỉ có tổng cộng $4$ chỗ ngồi.

\newpage
\problem{C1. Sắp Xếp Chỗ Ngồi (Phiên bản Dễ)}{Codeforces - 2232C1}

\section*{Đề bài}
Đây là phiên bản Dễ của bài toán. Điểm khác biệt giữa các phiên bản là ở phiên bản này, giới hạn của $n$, $x$, $s$, $t$ nhỏ hơn. Bạn chỉ có thể hack nếu bạn đã giải được tất cả các phiên bản của bài toán này.

Các bạn của Alice đã đến bữa tiệc, và bây giờ họ đang xếp hàng để vào dự tiệc.

Có $x$ cái bàn ở bữa tiệc, mỗi bàn có $s$ chỗ ngồi. Mỗi chỗ ngồi chỉ có thể chứa một người.

Mỗi người bạn có một trong ba tính cách sau:
\begin{itemize}
    \item Người hướng nội (\texttt{I}) phải ngồi vào một chiếc bàn trống.
    \item Người hướng ngoại (\texttt{E}) phải ngồi vào một chiếc bàn không trống.
    \item Người hướng trung (\texttt{A}) có thể ngồi ở bất kỳ bàn nào.
\end{itemize}

Ban đầu, mọi chỗ ngồi đều trống. Tuy nhiên, vì Alice đang ăn bánh, các bạn của cô ấy đã xếp thành một hàng, và Alice không thể thay đổi thứ tự của họ. Đối với mỗi người trong hàng, Alice phải xếp họ vào một cái bàn hoặc đuổi họ ra khỏi bữa tiệc. Mỗi người được xếp chỗ trước khi người tiếp theo được phân bàn.

Muốn bữa tiệc thật vui vẻ, Alice cần xếp chỗ cho càng nhiều người càng tốt. Hãy giúp cô ấy tìm số lượng bạn bè tối đa mà cô ấy có thể xếp chỗ trong bữa tiệc.

Lưu ý rằng một khi một người bạn đã ngồi xuống, họ không được phép di chuyển ngay cả khi chỗ ngồi hiện tại không còn phù hợp với tính cách của họ nữa.

\section*{Dữ liệu vào}
\begin{inputbox}
Mỗi test bao gồm nhiều test case. Dòng đầu tiên chứa số lượng test case $t$. Mô tả của các test case như sau.

Dòng đầu tiên của mỗi test case chứa ba số nguyên $n$, $x$, và $s$ – số lượng bạn bè của Alice, số lượng bàn, và số chỗ ngồi trên mỗi bàn.

Dòng thứ hai chứa một xâu $u$ có độ dài $n$ chỉ bao gồm các chữ cái \texttt{A}, \texttt{E}, và \texttt{I}, tương ứng đại diện cho người hướng trung, người hướng ngoại và người hướng nội.
\end{inputbox}

\section*{Dữ liệu ra}
\begin{outputbox}
Đối với mỗi test case, in ra một số nguyên: số người tối đa được xếp chỗ.
\end{outputbox}

\section*{Giới hạn}
\begin{constraintbox}
\begin{itemize}
    \item $1 \le t \le 500$
    \item $1 \le n, x, s \le 3000$
    \item Tổng của $n$ trong tất cả các test case không vượt quá $3000$.
\end{itemize}
\end{constraintbox}

\section*{Ví dụ}
\textbf{Input:}
\begin{lstlisting}
6
5 2 2
EIAIE
20 5 5
AEIEEEEIEAAEIEEEEIEA
8 2 4
AAAAAIEE
8 4 2
AIEAEAAI
8 3 3
AIEAEAAI
4 2 2
IAEE
\end{lstlisting}

\textbf{Output:}
\begin{lstlisting}
4
20
7
7
7
4
\end{lstlisting}

\explanation
Trong test case đầu tiên, có $2$ bàn, mỗi bàn có $2$ chỗ ngồi. Dưới đây là một trong những cách để đạt được số người tối đa được xếp chỗ:

\begin{itemize}
    \item Người đầu tiên là người hướng ngoại (\texttt{E}). Vì tất cả các bàn đều trống, họ phải rời khỏi bữa tiệc.
    \item Người thứ hai là người hướng nội (\texttt{I}). Alice có thể xếp họ vào bàn đầu tiên, đang trống.
    \item Người thứ ba là người hướng trung (\texttt{A}). Alice có thể xếp họ vào bàn đầu tiên.
    \item Người thứ tư là người hướng nội (\texttt{I}). Alice có thể xếp họ vào bàn thứ hai, đang trống.
    \item Người thứ năm là người hướng ngoại (\texttt{E}). Alice có thể xếp họ vào bàn thứ hai, lúc này đã không còn trống.
\end{itemize}

Như vậy, có bốn người được xếp chỗ trong bữa tiệc. Đây là số lượng tối đa vì bữa tiệc chỉ có bốn chỗ ngồi.

\newpage
\problem{B. Cake Leveling}{Codeforces - 2232B}

\section*{Đề bài}
Alice đang chuẩn bị một chiếc bánh cho bữa tiệc của mình. Tuy nhiên, cô ấy đang vội nên lớp kem trên mặt bánh không đều. Để giải quyết nhanh vấn đề này, Alice sẽ đặt dao ở một độ cao \textbf{nguyên} nào đó rồi gạt kem từ trái sang phải để làm phẳng mặt kem.

Cụ thể, gọi $a_i$ là độ cao của kem ở vị trí thứ $i$. Giả sử Alice đặt dao ở một độ cao \textbf{nguyên} $h$. Nếu độ cao của kem ở vị trí $i$ lớn hơn $h$, phần kem thừa sẽ bị đẩy sang vị trí $i+1$. Phần kem thừa ở vị trí $n$ sẽ bị đẩy ra khỏi bánh hoàn toàn.

Alice và bạn bè rất thích kem trên bánh. Vì Alice có thể quyết định chỉ phục vụ một phần tiền tố của chiếc bánh thay vì cả chiếc bánh, hãy giúp cô ấy tìm độ cao tối đa của kem có thể đạt được sao cho mặt kem phẳng đối với $i$ vị trí đầu tiên, với mọi $i = 1, 2, \ldots, n$.

\section*{Dữ liệu}
Mỗi test gồm nhiều test cases. Dòng đầu tiên chứa số lượng test cases $t$ ($1 \le t \le 10^4$). Mô tả của các test cases như sau.

Dòng đầu tiên của mỗi test case chứa một số nguyên $n$ ($2 \le n \le 2\cdot 10^5$) — độ dài của chiếc bánh của Alice.

Dòng thứ hai của mỗi test case chứa $n$ số nguyên $a_1, a_2, \ldots, a_n$ ($1 \le a_i \le 10^9$) — độ cao của kem ở vị trí thứ $i$.

Tổng của $n$ trên tất cả các test cases không vượt quá $2 \cdot 10^5$.

\section*{Kết quả}
Với mỗi test case, in ra $n$ số nguyên, trong đó số thứ $i$ thể hiện độ cao tối đa của kem có thể đạt được sao cho mặt kem phẳng đối với $i$ vị trí đầu tiên.

\section*{Ví dụ}

\begin{inputbox}5
3
4 2 3
5
2 3 4 3 2
5
3 3 3 1 1
3
913764826 346182673 764382516
8
6 7 6 7 6 7 6 7
\end{inputbox}\begin{outputbox}4 3 3 
2 2 2 2 2 
3 3 3 2 2 
913764826 629973749 629973749 
6 6 6 6 6 6 6 6 
\end{outputbox}%


\section*{Giải thích}
Giải thích cho test case thứ nhất:
\begin{itemize}
    \item Khi $i = 1$, Alice quan tâm đến phần bánh được biểu diễn bằng mảng $[4]$. Vì bánh đã phẳng sẵn, độ cao tối đa của kem là $4$.
    \item Khi $i = 2$, Alice quan tâm đến phần bánh được biểu diễn bằng mảng $[4, 2]$. Nếu Alice đặt dao ở độ cao $4$, bánh thu được sẽ có độ cao kem là $[4, 2]$, làm cho bánh không phẳng. Tuy nhiên, nếu Alice đặt dao ở độ cao $3$, một đơn vị kem sẽ bị đẩy từ vị trí thứ nhất sang vị trí thứ hai, làm cho bánh thu được có độ cao kem là $[3, 3]$, tức là đã phẳng. Do đó, độ cao tối đa của kem trong khi vẫn giữ phẳng là $3$.
    \item Khi $i = 3$, nếu Alice đặt dao ở độ cao $4$, độ cao kem thu được sẽ là $[4, 2, 3]$. Tuy nhiên, nếu Alice đặt dao ở độ cao $3$, độ cao kem thu được sẽ là $[3, 3, 3]$. Do đó, độ cao tối đa của kem trong khi vẫn giữ phẳng là $3$.
\end{itemize}

\newpage
\problem{Convergence}{Codeforces - 2232A}

\section*{Đề bài}
Alice đang mời bạn bè đến một bữa tiệc để ăn bánh. Tuy nhiên, mỗi người bạn có thể không ở cùng một địa điểm, do đó mọi người cần phải tập trung tại cùng một địa điểm trước.

Alice có $n$ người bạn, trong đó người bạn thứ $i$ đang ở vị trí $a_i$. Để tập hợp mọi người về cùng một chỗ, Alice phải thực hiện nhiều cuộc gọi nhóm. Thật không may, tín hiệu điện thoại rất yếu, do đó Alice chỉ có thể gọi cho $2$ người khác cùng một lúc.

Là một người tốt, Alice không muốn bạn bè của mình phải đi bộ quá xa. Vì vậy, với mỗi cuộc gọi nhóm gồm người bạn thứ $i$ và người bạn thứ $j$, Alice sẽ bảo cả hai người gặp nhau tại một vị trí nguyên nào đó nằm trong khoảng từ $\min(a_i, a_j)$ đến $\max(a_i, a_j)$ (bao gồm cả hai đầu mút). Sau đó, cả hai sẽ di chuyển đến vị trí đó với tốc độ cực nhanh, đến mức Alice không thể thực hiện bất kỳ cuộc gọi nhóm nào trong lúc họ đang di chuyển. Lưu ý rằng Alice có thể gọi lại cho những người bạn này sau khi họ đã đến vị trí đó.

Bữa tiệc sắp bắt đầu, nên Alice cần thực hiện các cuộc gọi nhóm thật nhanh. Hãy giúp cô ấy tìm số lượng cuộc gọi nhóm tối thiểu cần thực hiện.

\section*{Dữ liệu vào}
Mỗi test gồm nhiều test cases. Dòng đầu tiên chứa số lượng test cases $t$ ($1 \le t \le 500$). Phần mô tả các test cases theo sau.

Dòng đầu tiên của mỗi test case chứa một số nguyên $n$ ($2 \le n \le 100$) — số lượng bạn bè của Alice.

Dòng thứ hai của mỗi test case chứa $n$ số nguyên $a_1, a_2, \ldots, a_n$ ($1 \le a_i \le 10^9$) — vị trí của những người bạn.

\section*{Dữ liệu ra}
Với mỗi test case, in ra số lượng cuộc gọi nhóm tối thiểu mà Alice cần thực hiện để tất cả bạn bè của cô ấy tập trung tại cùng một địa điểm.

\section*{Ví dụ}


\begin{inputbox}4
5
1 2 3 4 5
5
1 1 1 2 2
11
3 1 4 1 5 9 2 6 5 3 5
5
1 2 2 2 2
\end{inputbox}
\begin{outputbox}2
2
5
1
\end{outputbox}



\section*{Giải thích}
Trong test case đầu tiên, số lượng cuộc gọi nhóm tối thiểu là $2$. Một cách để tập hợp mọi người tại cùng một vị trí như sau:
\begin{itemize}
    \item Gọi cho người bạn thứ $1$ và thứ $4$ và bảo họ di chuyển đến vị trí $3$. Vị trí của họ bây giờ là $[3, 2, 3, 3, 5]$.
    \item Gọi cho người bạn thứ $2$ và thứ $5$ và bảo họ di chuyển đến vị trí $3$. Vị trí của họ bây giờ là $[3, 3, 3, 3, 3]$.
\end{itemize}

Trong test case thứ hai, số lượng cuộc gọi nhóm tối thiểu là $2$. Một cách để tập hợp mọi người tại cùng một vị trí như sau:
\begin{itemize}
    \item Gọi cho người bạn thứ $1$ và thứ $4$ và bảo họ di chuyển đến vị trí $1$. Vị trí của họ bây giờ là $[1, 1, 1, 1, 2]$.
    \item Gọi cho người bạn thứ $4$ và thứ $5$ và bảo họ di chuyển đến vị trí $1$. Vị trí của họ bây giờ là $[1, 1, 1, 1, 1]$.
\end{itemize}

\newpage
\problem{F. Quadratic Jumps}{Codeforces - 2231F}

\section*{Đề bài}
Cho hai số nguyên $n$ và $q$. Xét một đồ thị có $n$ đỉnh, trong đó đỉnh $i$ và $j$ được nối với nhau bằng một cạnh khi và chỉ khi khoảng cách giữa chúng $|j - i|$ là một số chính phương$^*$.

Bạn được cho $q$ cặp số $a_i, b_i$. Với mỗi cặp trong số $q$ cặp này, bạn cần tìm khoảng cách ngắn nhất giữa đỉnh $a_i$ và $b_i$ trong đồ thị này. Có thể chứng minh được rằng đồ thị luôn liên thông, do đó khoảng cách giữa $a_i$ và $b_i$ không bao giờ là vô hạn.

$^*$Một số nguyên $x$ được gọi là số chính phương nếu tồn tại một số nguyên $y$ sao cho $x = y^2$.

\section*{Dữ liệu vào}
Mỗi test gồm nhiều test case. Dòng đầu tiên chứa số lượng test case $t$ ($1 \le t \le 1000$). Tiếp theo là phần mô tả của các test case.

Dòng đầu tiên của mỗi test case chứa hai số nguyên $n$ và $q$ ($2 \le n \le 2 \cdot 10^5$, $1 \le q \le 10^5$) — số đỉnh của đồ thị và số lượng cặp đỉnh cần tìm khoảng cách.

Sau đó, $q$ dòng tiếp theo mô tả các cặp đỉnh cần tìm khoảng cách ngắn nhất. Mỗi cặp được biểu diễn bởi hai số $a, b$ ($1 \leq a < b \leq n$) — số thứ tự của các đỉnh cần tính khoảng cách ngắn nhất giữa chúng.

\begin{constraintbox}
Tổng của $n$ qua tất cả các test case không vượt quá $2 \cdot 10^5$, và tổng của $q$ qua tất cả các test case không vượt quá $10^5$.
\end{constraintbox}

\section*{Dữ liệu ra}
Với mỗi test case, in ra khoảng cách ngắn nhất cho từng cặp đỉnh trong số $q$ truy vấn trên các dòng riêng biệt.

\section*{Ví dụ}
\begin{inputbox}
2
5 4
1 2
1 3
1 4
1 5
8 3
3 7
2 5
1 7
\end{inputbox}
\begin{outputbox}
1
2
2
1
1
2
3
\end{outputbox}

\explanation

\textbf{Test case 1}: Với $n = 5$, các cạnh được tạo ra giữa các đỉnh có chênh lệch là số chính phương (1 hoặc 4):
\begin{itemize}
    \item Chênh lệch 1: $(1, 2), (2, 3), (3, 4), (4, 5)$.
    \item Chênh lệch 4: $(1, 5)$.
\end{itemize}
Xét $q=4$ truy vấn:
\begin{enumerate}
    \item $1 \to 2$: Có cạnh nối trực tiếp vì $|2 - 1| = 1 = 1^2$. Khoảng cách là 1.
    \item $1 \to 3$: Đi qua đường đi $1 \to 2 \to 3$. Khoảng cách là 2.
    \item $1 \to 4$: Đi qua đường đi $1 \to 5 \to 4$ (vì $1 \to 5$ có khoảng cách là 4, $5 \to 4$ có khoảng cách là 1). Khoảng cách là 2.
    \item $1 \to 5$: Có cạnh nối trực tiếp vì $|5 - 1| = 4 = 2^2$. Khoảng cách là 1.
\end{enumerate}

\textbf{Test case 2}: Với $n = 8$, các cạnh được tạo ra giữa các đỉnh có chênh lệch là 1 hoặc 4:
\begin{itemize}
    \item Chênh lệch 1: $(1, 2), (2, 3), (3, 4), (4, 5), (5, 6), (6, 7), (7, 8)$.
    \item Chênh lệch 4: $(1, 5), (2, 6), (3, 7), (4, 8)$.
\end{itemize}
Xét $q=3$ truy vấn:
\begin{enumerate}
    \item $3 \to 7$: Có cạnh nối trực tiếp vì $|7 - 3| = 4 = 2^2$. Khoảng cách là 1.
    \item $2 \to 5$: Đường đi ngắn nhất là $2 \to 6 \to 5$ hoặc $2 \to 1 \to 5$. Khoảng cách là 2.
    \item $1 \to 7$: Đường đi ngắn nhất có thể là $1 \to 2 \to 6 \to 7$. Khoảng cách là 3.
\end{enumerate}

\newpage


\end{document}
